\section{Experiments}
\label{sec:experiments}

% ---------------------------------------------------------------------------
% NOTHING IN THIS SECTION MAY CONTAIN A NUMBER THAT AN EXPERIMENT DID NOT
% PRODUCE. Leave \ARnum{} and the commented-out environments standing until a
% run has actually finished and its log is committed alongside the paper. A
% fabricated table is the one failure mode this pipeline cannot recover from.
% ---------------------------------------------------------------------------

\subsection{Setup}
\label{sec:setup}

\paragraph{Models and data.}
\ARTODO{exact checkpoints, datasets, splits and sizes}

\paragraph{Baselines.}
% At least one baseline must be the obvious thing a skeptical reviewer would
% try instead of the proposed method.
\ARTODO{baselines and why these are the right comparisons}

\paragraph{Metrics.}
\ARTODO{metrics, how they are computed, and what counts as a meaningful gap}

\paragraph{Implementation.}
\ARTODO{hardware, precision, seeds, hyperparameters, and the number of repeats}

\subsection{Main Results}
\label{sec:main-results}

\ARTODO{one paragraph reading the table out loud: what moved, by how much, and where it did not}

% Uncomment and fill once the run has finished. Keep the caption self-contained:
% a reader should understand the table without returning to the body text.
%
% \begin{table}[t]
%   \centering
%   \caption{Main results. Mean over N seeds, standard deviation in parentheses.}
%   \label{tab:main}
%   \begin{tabular}{lcc}
%     \toprule
%     Method & Metric A $\uparrow$ & Metric B $\downarrow$ \\
%     \midrule
%     Baseline &  &  \\
%     Ours     &  &  \\
%     \bottomrule
%   \end{tabular}
% \end{table}

\ARfig{Main result figure: the headline comparison. Produce this from a script committed in the repository so it can be regenerated.}

\subsection{Ablations}
\label{sec:ablations}

% One ablation per design decision claimed as a contribution. If a component
% cannot be ablated away, explain why rather than omitting it.
\ARTODO{ablation table and the conclusion each row supports}

% \begin{table}[t]
%   \centering
%   \caption{Ablations. Each row removes one component from the full method.}
%   \label{tab:ablation}
%   \begin{tabular}{lc}
%     \toprule
%     Variant & Metric A $\uparrow$ \\
%     \midrule
%     Full method &  \\
%     $-$ component 1 &  \\
%     $-$ component 2 &  \\
%     \bottomrule
%   \end{tabular}
% \end{table}

\subsection{Analysis}
\label{sec:analysis}

% The section that separates a strong paper from a results dump: why does the
% method behave the way the tables show? Include the case where it does not.
\ARTODO{mechanism evidence, scaling behaviour, and failure cases}

\ARfig{Analysis figure: the evidence for the mechanism claimed in Section~\ref{sec:intuition}.}
